\documentclass{article}
\usepackage[utf8]{inputenc}
\usepackage[T2A]{fontenc}
\usepackage[utf8]{inputenc}
\usepackage[english,russian]{babel}
\usepackage[a4paper,top=2.0cm,bottom=2.0cm,left=3.0cm,right=1.5cm]{geometry}
\usepackage{indentfirst}
\usepackage{titlesec}
\usepackage{amsmath}
\usepackage{esvect}
\usepackage{subfig}
\usepackage{listings}
\usepackage{subcaption}
\usepackage{color,soul}
\usepackage{graphicx}
\usepackage{float}
\usepackage{cmap}
\usepackage{mathtext} 
\usepackage{lastpage}
\usepackage{totcount}
\usepackage{fancyhdr}
\captionsetup{labelsep=endash}
\captionsetup[figure]{name=Рисунок}
\captionsetup{font={large}}

% \usepackage[colorlinks=true, allcolors=blue]{hyperref}
\usepackage{float}
\renewcommand{\floatpagefraction}{1}

\fancypagestyle{plain}{%
\fancyhf{} % clear all header and footer fields
\fancyfoot[C]{\fontsize{14pt}{14pt}\selectfont\thepage} % except the center
\renewcommand{\headrulewidth}{0pt}
\renewcommand{\footrulewidth}{0pt}}
\pagestyle{plain}

\renewcommand{\thesubfigure}{\asbuk{subfigure}}
\renewcommand\thefigure{\arabic{section}.\arabic{figure}}
\makeatletter
\@addtoreset{figure}{section} % Нумерация рисунков будет пересчитываться при каждом новом разделе
\makeatother

\linespread{1.5}
\setlength{\parskip}{0.4cm}
\setlength{\parindent}{1cm}

\begin{document}

\newcommand{\refpar}[1]{(\ref{#1})}
\fontsize{14pt}{14pt}\normalfont

\titleformat{\subsection}
{\normalfont\fontsize{14pt}{14pt}\selectfont\bfseries}{\thesubsection}{1em}{}
\titleformat{\subsubsection}
{\normalfont\fontsize{14pt}{14pt}\selectfont\bfseries}{\thesubsubsection}{1em}{}

\begin{titlepage}
{
\fontsize{14pt}{20pt}\normalfont
\centerline{МИНИСТЕРСТВО НАУКИ И ВЫСШЕГО ОБРАЗОВАНИЯ РОССИЙСКОЙ}
\centerline{ФЕДЕРАЦИИ}
\centerline{федеральное государственное автономное образовательное учреждение}
\centerline{высшего образования}
\centerline{УРАЛЬСКИЙ ФЕДЕРАЛЬНЫЙ УНИВЕРСИТЕТ}
\centerline{имени первого Президента России Б.Н. Ельцина}
\vskip0.5cm
\centerline{ИНСТИТУТ ЕСТЕСТВЕННЫХ НАУК И МАТЕМАТИКИ}
\vskip1cm
\centerline{\hl{Кафедра теоретической и математической физики}}
}
\vskip1cm

\centerline{\textbf{Поиск оптимальной стратегии в симуляторе владельца продукта}}
\centerline{\textbf{с помощью Deep Reinforcement Learning}}

\vskip0.5cm
\centerline{Направление подготовки 02.04.01 <<Математика и компьютерные науки>>}

\null\hfill

\noindent\makebox[\textwidth][c]{%
    \begin{minipage}[t]{.5\textwidth}
        \centerline{Зав. кафедрой:}
        \centerline{\hl{д. ф.-м. н., доц. Е. А. Елфимова}}
        \vskip0.8cm
        \centerline{\rule[-1pt]{4.5cm}{0.4pt}}
    \end{minipage}%
    
    \begin{minipage}[t]{.5\textwidth}
        \centerline{Выпускная квалификационная}
        \centerline{работа магистранта}
        \centerline{\textbf{Крутовского}}
        \centerline{\textbf{Данила Александровича}}
        \vskip0.4cm
        \centerline{\rule[-1pt]{4.5cm}{0.4pt} }
    \end{minipage}
}

\vskip0.5cm
\noindent\makebox[\textwidth][c]{%
    \begin{minipage}[t]{.5\textwidth}
        \centerline{Нормоконтролер:}
        \centerline{\hl{к. ф.-м. н. Е. С. Пьянзина}}
        \vskip0.8cm
        \centerline{\rule[-1pt]{4.5cm}{0.4pt}}
    \end{minipage}% <---------------- Note the use of "%"
    
    \begin{minipage}[t]{.5\textwidth}
        \centerline{Научный руководитель:\\}
        \centerline{к. ф.-м. н. А. А. Кошелев\\}
        \vskip0.8cm
        \centerline{\rule[-1pt]{4.5cm}{0.4pt}} 
    \end{minipage}
}

\vfill
\centerline{Екатеринбург}
\centerline{2025}
\end{titlepage}

\setcounter{page}{2}

% \regtotcounter{figure}
\regtotcounter{table}
\newtotcounter{citnum}
\def\oldbibitem{} 
\let\oldbibitem=\bibitem
\def\bibitem{\stepcounter{citnum}\oldbibitem}

\section*{\centering{РЕФЕРАТ}}
Выпускная квалификационная работа Крутовского Данил Александровича «Поиск оптимальной стратегии в симуляторе владельца продукта с помощью Deep Reinforcement Learning». В работе содержится \pageref{LastPage} страниц, 28 рисунков, \total{table} таблица, \total{citnum} источников, 1 приложение.

Ключевые слова: обучение с подкреплением, среда, агент, функция наград, Q-обучение.

Целью данной работы является алгоритм, способный обыграть человека в Симуляторе владельца продукта.

\hl{Написать реферат}

Для достижения цели были поставлены следующие задачи:
\begin{enumerate}
  \item Изучение теоретического материала по темам: обучение с подкреплением;
  \item Настройка взаимодействие агента с игрой-симулятором владельца продукта;
  \item Построение бейзлайна обучения агента на основе Double DQN;
  \item Исследование проблем текущего пайплайна обучения агента.
\end{enumerate}

\newpage

\let\mtcontentsname\contentsname
\renewcommand\contentsname{\centering\MakeUppercase\mtcontentsname}
\tableofcontents
\newpage

\section*{\centering{ОБОЗНАЧЕНИЯ И СОКРАЩЕНИЯ}}
\addcontentsline{toc}{section}{ОБОЗНАЧЕНИЯ И СОКРАЩЕНИЯ}

В настоящей работе применяют следующие обозначения и сокращения:

RL -- reinforcement learning, обучение с подкрепление

$\pi$ -- стратегия, политика 

DQN - Deep Q-Network, алгоритм обучения с подкреплением

Python - 

pytorch - 

Godot - 

Датасет - 

Лидерборд - 



\hl{Дополнить списко обозначений}

Репозиторий -- хранилище файлов, кода или данных, где разработчики могут совместно работать над проектом, отслеживать изменения, управлять версиями и делиться информацией

GitHub -- веб-платформа для хостинга репозиториев и совместной работы над проектами с использованием системы контроля версий Git

Git -- распределенная система управления версиями, предназначенная для отслеживания изменений в исходном коде проекта, совместной работы нескольких разработчиков и управления версиями файлов

Юнит-тесты -- процесс проверки отдельных модулей программы на корректность их работы

\newpage

\section*{\centering{ВВЕДЕНИЕ}}
\addcontentsline{toc}{section}{ВВЕДЕНИЕ}

Обучение с подкреплением - это обучение тому, что надо делать, как следует отображать ситуацию в действия, чтобы максимизировать некоторый сигнал поощрения (вознаграждения), принимающий числовые значения \cite{Satton-Barto}.

Product Owner Sim \cite{game-link} предлагает вам управлять проектом с бюджетом 200 тысяч долларов и двумя роботами-разработчиками.
Каждый разработчик тратит на задачи по 10 часов за спринт, но при условии, что они требуют оплаты в 10 тысяч долларов за работу.
На доход проекта влияет количество пользователей и их лояльность.
Для увеличения этих показателей нужно покупать исследования, декомпозировать их на задачи и выпускать новые функции.
Вы должны использовать свои ресурсы оптимально для достижения цели в один миллион долларов.
Главный показатель в игре -- номер спринта, на котором была достигнута цель, чем меньше, тем лучше.

\hl{Дополнить или переписать введение}

\newpage

\section*{\centering{ОСНОВНАЯ ЧАСТЬ}}
\addcontentsline{toc}{section}{ОСНОВНАЯ ЧАСТЬ}

\section{Подготовка инфраструктуры}

В настоящее время для решения задач машинного обучения широко используются инструменты на основе языка программирования Python, в частности библиотека PyTorch.

Особенностью задачи обучения с подкреплением является отсутствие классических датасетов. 
Агент обучается посредством взаимодействия со средой, поэтому ключевым аспектом является настройка эффективного механизма получения данных о состоянии среды. 
Это необходимо для обеспечения возможности быстрого обучения моделей и валидации гипотез в разумные временные рамки.

\subsection{Взаимодействие с игрой}

В нашем случае среда представлена игрой, разработанной на движке Godot. 
У нас был доступ к исходному коду игры, что позволило провести более глубокую интеграцию.

Для обеспечения максимально эффективного взаимодействия модели со средой было принято решение переписать логику игры на языке Python. 
Это позволило избежать задержек, связанных с долгим отображением анимаций в игре, и дало возможность более детально изучить реализацию игровых механик. 
В результате скорость применения действий агентом в среде увеличилась в тысячи раз.

Альтернативные подходы к взаимодействию с игрой могли включать:

1. Симуляция взаимодействия через пользовательский интерфейс. 
Этот метод предполагает имитацию действий пользователя, таких как перемещения курсора и нажатия мыши, для взаимодействия с элементами игры. 
Такой подход мог бы быть полезен в дальнейшем, если бы возникла необходимость интегрировать агента в официальный рейтинг игроков (лидерборд) игры, так как он позволяет модели взаимодействовать с игрой так же, как это делает человек.

2. Взаимодействие с игрой на уровне бинарных данных. 
Имея доступ к исходному коду игры, можно было бы реализовать взаимодействие на более глубоком уровне — путём сбора и обработки бинарных данных. 
Это позволило бы модели работать с игрой как с библиотечной функцией, что потенциально может ускорить обмен данными и упростить процесс обучения агента, минуя необходимость в визуальной интерпретации состояния игры.

\subsection{Минимальный подходящий алгоритм}

\hl{написать содержание магистерской}

\section{Поиск системы наград}

\subsection{Наказывающая система наград}
\subsection{Поощряющая система наград}
\subsection{Дальновидная система наград}

\section{Оптимизация процесса обучения модели}
\subsection{Маскирование недоступных действий}
\subsection{Сбора задач в спринте с помощью алгоритма "Задача о рюкзаке"}
\subsection{Добавления действия, позволяющего доигрывать пустыми спринтами}

\section{Сравнение с человеческими возможностями}
\subsection{Взаимодействие с веб версией}
\subsection{Подготовка фиксированной среды}

\newpage
\section*{\centering{ЗАКЛЮЧЕНИЕ}}
\addcontentsline{toc}{section}{ЗАКЛЮЧЕНИЕ}

\hl{Написать заключение}

\newpage

\renewcommand\refname{\centering{СПИСОК ИСПОЛЬЗОВАННЫХ ИСТОЧНИКОВ И ЛИТЕРАТУРЫ}}
\addcontentsline{toc}{section}{СПИСОК ИСПОЛЬЗОВАННЫХ ИСТОЧНИКОВ И ЛИТЕРАТУРЫ}

\bibliographystyle{unsrt}
\bibliography{references} 

\newpage
\section*{\centering{ПРИЛОЖЕНИЕ А}}
\addcontentsline{toc}{section}{ПРИЛОЖЕНИЕ А}
\lstset{
  language=Python,
  basicstyle=\fontsize{14}{16}\linespread{1}\selectfont\ttfamily,
  basewidth=0.5em,
  aboveskip=0pt, belowskip=0pt, showstringspaces=false,
}
\begin{lstlisting}
import datetime
import os
import sys

import pandas as pd

from typing import List

from pipeline.study_agent import save_dqn_agent

sys.path.append("..")
sys.path.append("../..")

from algorithms import DoubleDQN
from environment import CreditPayerEnv
from environment.backlog_env import BacklogEnv
from environment.userstory_env import UserstoryEnv
from environment.reward_sytem import (
    EmpiricalCreditStageRewardSystem,
    FullPotentialCreditRewardSystem,
)
from pipeline.aggregator_study import update_reward_system_config
from pipeline import LoggingStudy


def make_credit_study(trajectory_max_len, episode_n, potential):
    userstory_env = UserstoryEnv(4, 0, 0)
    backlog_env = BacklogEnv(12, 0, 0, 0, 0, 0)
    if potential:
        reward_system = FullPotentialCreditRewardSystem(config={}, coefficient=1)
    else:
        reward_system = EmpiricalCreditStageRewardSystem(True, config={})
    env = CreditPayerEnv(
        userstory_env,
        backlog_env,
        with_end=True,
        with_info=True,
        reward_system=reward_system,
    )
    update_reward_system_config(env, reward_system)

    state_dim = env.state_dim
    action_n = env.action_n

    agent = DoubleDQN(
        state_dim,
        action_n,
        gamma=0.9,
        tau=0.001,
        epsilon_decrease=1e-4,
        batch_size=64,
        lr=1e-3,
        epsilon_min=0.01,
    )

    study = LoggingStudy(env, agent, trajectory_max_len)
    study.study_agent(episode_n)

    return study


def main(potential):
    episode_n = 1501
    study = make_credit_study(200, episode_n, potential)
    save_dqn_agent(study.agent, "models/credit_start_model.pt")


if __name__ == "__main__":
    n = 1
    for i in range(n):
        main(True)

\end{lstlisting}

\end{document}
